%% =====================================================================
%%  T-14-04  The Four Amplifiers
%%  -------------------------------------------------------------------
%%  One idea per file. Core is mandatory. Slots in canonical order.
%%  Max 700 words. No layout commands. No filler. New source material
%%  for this entry goes in sources/<BOOK>/T-14-04.tex -- never by
%%  rewriting this file (docs/RULES.md R0.1).
%% =====================================================================
%% @kind: topic
%% @id: T-14-04
%% @title: The Four Amplifiers
%% @subtitle: Written, active, public and effortful commitments bind hardest
%% @chapter: c14
%% @order: 40
%% @status: stable
%% @sources: B4
%% @updated: 2026-09-19
%% @allow-rewrite: slot migration to the seven-question format (docs/RULES.md section 2)

\topic{T-14-04}{The Four Amplifiers}{Written, active, public and effortful commitments bind hardest}


\begin{core}
Not all commitments hold equally. Four properties strengthen one: it is written down, it was performed rather than assented to, other people saw it, and it cost something to make. Each can be designed for, and each makes the commitment harder to abandon than its content warrants.
\end{core}

\begin{mechanism}
Writing supplies a record the person cannot later re-describe, and the act of writing is itself a small piece of authorship --- the position stops being something heard and becomes something produced. Action binds more than assent because a performed behaviour is evidence about the self, whereas agreement is only evidence about a conversation. Publicity adds an audience whose expectations are now part of the cost of reversal, and makes the commitment a statement of identity rather than of preference. Effort is the strongest of the four because the person must justify the cost they paid, and the cheapest justification available is that they wanted to pay it.

The four compound. A commitment that is all four is close to unbreakable from inside, which is why institutions that need reliable allegiance arrange all four rather than any one.
\end{mechanism}

\begin{application}
  \item Get it in the counterpart's own handwriting wherever possible, even for something trivial.
  \item Arrange for them to act before they decide, so the decision follows the behaviour.
  \item Introduce an audience early, and let the counterpart make the announcement themselves.
  \item Invest their effort visibly. Effort they can see is effort they must account for.
\end{application}

\begin{feedback}
  \item You are asked to sign, initial, fill in or write out the agreement yourself rather than receive it complete.
  \item You perform something before you have decided --- a trial, a configuration, an assessment, a visit.
  \item The commitment is announced to others, or others are invited to witness it.
  \item Considerable time, money or effort has already been spent, and leaving would forfeit it.
  \item Refusing now would require explaining yourself to people who already know you agreed.
\end{feedback}

\begin{failure}
Amplifying a commitment you do not fully control is dangerous in both directions. The same properties that bind a counterpart bind you: a public written position of your own is expensive to revise, and people who design these structures usually find themselves inside one. There is also a straightforward honesty cost --- a commitment extracted by engineering its amplifiers is not the counterpart's commitment, and discovering that is the end of the relationship rather than of the deal.
\end{failure}

\begin{countermeasures}
  \item Decline to write, sign or perform anything before the decision is made. Ask for the terms complete, in their hand, not yours.
  \item Keep commitments private until they are settled. An audience converts a preference into an identity claim.
  \item Treat sunk effort as spent rather than as a reason. Ask what you would do if the effort had never been spent.
  \item Reverse in writing. A written commitment is answered by a written withdrawal, which restores symmetry and removes the record's advantage.
\end{countermeasures}

\sources{B4}
\seesources
\seealso{T-14-03, T-06-04, T-24-02, T-25-01}
